% Options for packages loaded elsewhere
% Options for packages loaded elsewhere
\PassOptionsToPackage{unicode}{hyperref}
\PassOptionsToPackage{hyphens}{url}
%
\documentclass[
  11pt,
  letterpaper,
  DIV=11,
  numbers=noendperiod,
  twoside]{scrartcl}
\usepackage{xcolor}
\usepackage[left=1.1in, right=1in, top=0.8in, bottom=0.8in,
paperheight=9.5in, paperwidth=7in, includemp=TRUE, marginparwidth=0in,
marginparsep=0in]{geometry}
\usepackage{amsmath,amssymb}
\setcounter{secnumdepth}{3}
\usepackage{iftex}
\ifPDFTeX
  \usepackage[T1]{fontenc}
  \usepackage[utf8]{inputenc}
  \usepackage{textcomp} % provide euro and other symbols
\else % if luatex or xetex
  \usepackage{unicode-math} % this also loads fontspec
  \defaultfontfeatures{Scale=MatchLowercase}
  \defaultfontfeatures[\rmfamily]{Ligatures=TeX,Scale=1}
\fi
\usepackage{lmodern}
\ifPDFTeX\else
  % xetex/luatex font selection
  \setmathfont[]{Garamond-Math}
\fi
% Use upquote if available, for straight quotes in verbatim environments
\IfFileExists{upquote.sty}{\usepackage{upquote}}{}
\IfFileExists{microtype.sty}{% use microtype if available
  \usepackage[]{microtype}
  \UseMicrotypeSet[protrusion]{basicmath} % disable protrusion for tt fonts
}{}
\usepackage{setspace}
% Make \paragraph and \subparagraph free-standing
\makeatletter
\ifx\paragraph\undefined\else
  \let\oldparagraph\paragraph
  \renewcommand{\paragraph}{
    \@ifstar
      \xxxParagraphStar
      \xxxParagraphNoStar
  }
  \newcommand{\xxxParagraphStar}[1]{\oldparagraph*{#1}\mbox{}}
  \newcommand{\xxxParagraphNoStar}[1]{\oldparagraph{#1}\mbox{}}
\fi
\ifx\subparagraph\undefined\else
  \let\oldsubparagraph\subparagraph
  \renewcommand{\subparagraph}{
    \@ifstar
      \xxxSubParagraphStar
      \xxxSubParagraphNoStar
  }
  \newcommand{\xxxSubParagraphStar}[1]{\oldsubparagraph*{#1}\mbox{}}
  \newcommand{\xxxSubParagraphNoStar}[1]{\oldsubparagraph{#1}\mbox{}}
\fi
\makeatother


\usepackage{longtable,booktabs,array}
\newcounter{none} % for unnumbered tables
\usepackage{calc} % for calculating minipage widths
% Correct order of tables after \paragraph or \subparagraph
\usepackage{etoolbox}
\makeatletter
\patchcmd\longtable{\par}{\if@noskipsec\mbox{}\fi\par}{}{}
\makeatother
% Allow footnotes in longtable head/foot
\IfFileExists{footnotehyper.sty}{\usepackage{footnotehyper}}{\usepackage{footnote}}
\makesavenoteenv{longtable}
\usepackage{graphicx}
\makeatletter
\newsavebox\pandoc@box
\newcommand*\pandocbounded[1]{% scales image to fit in text height/width
  \sbox\pandoc@box{#1}%
  \Gscale@div\@tempa{\textheight}{\dimexpr\ht\pandoc@box+\dp\pandoc@box\relax}%
  \Gscale@div\@tempb{\linewidth}{\wd\pandoc@box}%
  \ifdim\@tempb\p@<\@tempa\p@\let\@tempa\@tempb\fi% select the smaller of both
  \ifdim\@tempa\p@<\p@\scalebox{\@tempa}{\usebox\pandoc@box}%
  \else\usebox{\pandoc@box}%
  \fi%
}
% Set default figure placement to htbp
\def\fps@figure{htbp}
\makeatother


% definitions for citeproc citations
\NewDocumentCommand\citeproctext{}{}
\NewDocumentCommand\citeproc{mm}{%
  \begingroup\def\citeproctext{#2}\cite{#1}\endgroup}
\makeatletter
 % allow citations to break across lines
 \let\@cite@ofmt\@firstofone
 % avoid brackets around text for \cite:
 \def\@biblabel#1{}
 \def\@cite#1#2{{#1\if@tempswa , #2\fi}}
\makeatother
\newlength{\cslhangindent}
\setlength{\cslhangindent}{1.5em}
\newlength{\csllabelwidth}
\setlength{\csllabelwidth}{3em}
\newenvironment{CSLReferences}[2] % #1 hanging-indent, #2 entry-spacing
 {\begin{list}{}{%
  \setlength{\itemindent}{0pt}
  \setlength{\leftmargin}{0pt}
  \setlength{\parsep}{0pt}
  % turn on hanging indent if param 1 is 1
  \ifodd #1
   \setlength{\leftmargin}{\cslhangindent}
   \setlength{\itemindent}{-1\cslhangindent}
  \fi
  % set entry spacing
  \setlength{\itemsep}{#2\baselineskip}}}
 {\end{list}}
\usepackage{calc}
\newcommand{\CSLBlock}[1]{\hfill\break\parbox[t]{\linewidth}{\strut\ignorespaces#1\strut}}
\newcommand{\CSLLeftMargin}[1]{\parbox[t]{\csllabelwidth}{\strut#1\strut}}
\newcommand{\CSLRightInline}[1]{\parbox[t]{\linewidth - \csllabelwidth}{\strut#1\strut}}
\newcommand{\CSLIndent}[1]{\hspace{\cslhangindent}#1}



\setlength{\emergencystretch}{3em} % prevent overfull lines

\providecommand{\tightlist}{%
  \setlength{\itemsep}{0pt}\setlength{\parskip}{0pt}}



 


\setlength\heavyrulewidth{0ex}
\setlength\lightrulewidth{0ex}
\usepackage[automark]{scrlayer-scrpage}
\clearpairofpagestyles
\cehead{
  Brian Weatherson
  }
\cohead{
  Countable Additivity and Accuracy Dominance
  }
\ohead{\bfseries \pagemark}
\cfoot{}
\makeatletter
\newcommand*\NoIndentAfterEnv[1]{%
  \AfterEndEnvironment{#1}{\par\@afterindentfalse\@afterheading}}
\makeatother
\NoIndentAfterEnv{itemize}
\NoIndentAfterEnv{enumerate}
\NoIndentAfterEnv{description}
\NoIndentAfterEnv{quote}
\NoIndentAfterEnv{equation}
\NoIndentAfterEnv{longtable}
\NoIndentAfterEnv{abstract}
\renewenvironment{abstract}
 {\vspace{-1.25cm}
 \quotation\small\noindent\emph{Abstract}:}
 {\endquotation}
\newfontfamily\tfont{EB Garamond}
\addtokomafont{disposition}{\rmfamily}
\addtokomafont{descriptionlabel}{\rmfamily}
\addtokomafont{title}{\normalfont\itshape}
\let\footnoterule\relax

\makeatletter
\renewcommand{\@maketitle}{%
  \newpage
  \null
  \vskip 2em%
  \begin{center}%
  \let \footnote \thanks
    {\itshape\huge\@title \par}%
    \vskip 0.5em%  % Reduced from default
    {\large
      \lineskip 0.3em%  % Reduced from default 0.5em
      \begin{tabular}[t]{c}%
        \@author
      \end{tabular}\par}%
    \vskip 0.5em%  % Reduced from default
    {\large \@date}%
  \end{center}%
  \par
  }
\makeatother
\RequirePackage{lettrine}

\renewenvironment{abstract}
 {\quotation\small\noindent\emph{Abstract}:}
 {\endquotation\vspace{-0.02cm}}

\setmainfont{EB Garamond Math}[
  BoldFont = {EB Garamond SemiBold},
  ItalicFont = {EB Garamond Italic},
  RawFeature = {+smcp},
]

\newfontfamily\scfont{EB Garamond Regular}[RawFeature=+smcp]
\renewcommand{\textsc}[1]{{\scfont #1}}

\renewcommand{\LettrineTextFont}{\scfont}
% Number theorem-like environments straight through (Theorem 1, Theorem 2, ...)
% rather than by section (1.1, 1.2, ...), which is Quarto's LaTeX default.
% Guarded, since Quarto only defines a counter for environments the document uses.
\makeatletter
\AtBeginDocument{%
  \@for\bw@thm:={theorem,lemma,corollary,proposition,conjecture,definition,%
                 example,exercise,refremark,refsolution}\do{%
    \@ifundefined{c@\bw@thm}{}{\expandafter\counterwithout\expandafter{\bw@thm}{section}}%
  }%
}
\makeatother
\KOMAoption{captions}{tableheading}
\makeatletter
\@ifpackageloaded{caption}{}{\usepackage{caption}}
\AtBeginDocument{%
\ifdefined\contentsname
  \renewcommand*\contentsname{Table of contents}
\else
  \newcommand\contentsname{Table of contents}
\fi
\ifdefined\listfigurename
  \renewcommand*\listfigurename{List of Figures}
\else
  \newcommand\listfigurename{List of Figures}
\fi
\ifdefined\listtablename
  \renewcommand*\listtablename{List of Tables}
\else
  \newcommand\listtablename{List of Tables}
\fi
\ifdefined\figurename
  \renewcommand*\figurename{Figure}
\else
  \newcommand\figurename{Figure}
\fi
\ifdefined\tablename
  \renewcommand*\tablename{Table}
\else
  \newcommand\tablename{Table}
\fi
}
\@ifpackageloaded{float}{}{\usepackage{float}}
\floatstyle{ruled}
\@ifundefined{c@chapter}{\newfloat{codelisting}{h}{lop}}{\newfloat{codelisting}{h}{lop}[chapter]}
\floatname{codelisting}{Listing}
\newcommand*\listoflistings{\listof{codelisting}{List of Listings}}
\usepackage{amsthm}
\theoremstyle{plain}
\newtheorem{theorem}{Theorem}[section]
\theoremstyle{remark}
\AtBeginDocument{\renewcommand*{\proofname}{Proof}}
\newtheorem*{remark}{Remark}
\newtheorem*{solution}{Solution}
\newtheorem{refremark}{Remark}[section]
\newtheorem{refsolution}{Solution}[section]
\makeatother
\makeatletter
\makeatother
\makeatletter
\@ifpackageloaded{caption}{}{\usepackage{caption}}
\@ifpackageloaded{subcaption}{}{\usepackage{subcaption}}
\makeatother
\usepackage{bookmark}
\IfFileExists{xurl.sty}{\usepackage{xurl}}{} % add URL line breaks if available
\urlstyle{same}
% fallback for those not using the hyperref driver hyperxmp:
\makeatletter
\@ifundefined{xmpquote}{\newcommand{\xmpquote}[1]{#1}}{}
\makeatother
\hypersetup{
  pdftitle={Countable Additivity and Accuracy Dominance},
  hidelinks,
  pdfcreator={LaTeX via pandoc}}


\title{Countable Additivity and Accuracy Dominance}
\author{}
\date{}
\begin{document}
\maketitle


\setstretch{1.1}
\section{Introduction}\label{introduction}

Standard presentations of the accuracy dominance argument for
probabilism, going back to work by James Joyce
(\citeproc{ref-Joyce1998}{1998}), show that undominated credences must
be \emph{finitely} additive. Can we generalise them to show that
undominated credences must be \emph{countably} additive? The question
has some interest for two reasons. First, a number of questions in
formal epistemology turn on whether credences are countably or merely
finitely additive. For example, some uses of indifference reasoning are
ruled out if credences must be countably additive
(\citeproc{ref-Ross2010}{Ross 2010}). Second, while finite additivity
strikes many people as intuitively plausible, countable additivity is
much more contested. This is related to the first reason. Many people
think it is reasonable to be certain that one of a countable infinity of
states will come about, and regard each of those states as equally
likely. In short, there can be a fair lottery over the naturals. That is
coherent only if countable additivity is not a constraint.

The obvious way to generalise Joyce's argument to countable additivity
would use a measure. The theorist comes up with some inaccuracy measure
\(\mathcal{I}(c, w)\), meaning the inaccuracy of credal state \(c\) in
world \(w\), and shows that if \(c\) is not countably additive, then
there is some alternative credal state \(c^\prime\) such that
\(\mathcal{I}(c, w) > \mathcal{I}(c^\prime, w)\) for all \(w\). This
approach has not been particularly successful.\footnote{This needs
  filling in. The paper should engage with Kelley
  (\citeproc{ref-Kelley2023}{2023}), Kelley
  (\citeproc{ref-Kelley2026}{2026}), Nielsen
  (\citeproc{ref-Nielsen2022}{2022}), Nielsen
  (\citeproc{ref-Nielsen2023infinite}{2023}), Pruss
  (\citeproc{ref-Pruss2022}{2022}), Kelley and Neth
  (\citeproc{ref-KelleyNeth2023}{2023}) and Park and Jung
  (\citeproc{ref-ParkJung2023}{2023}).}

The accuracy dominance argument does not, however, need a numerical
accuracy measure. All that's needed is a partial orderings of credal
states with respect to accuracy at each world. (Equivalently, a function
from worlds to partial orderings.) Instead of looking for a function
\(\mathcal{I}(c, w)\), we can look for a relation
\(\succ_{\mathcal{I}, w}\), where \(c \succ_{\mathcal{I}, w} c^\prime\)
should be read as saying that \(c\) is more inaccuracte than
\(c^\prime\) at \(w\). As long as we can (a) justify that
\(\succ_{\mathcal{I}, w}\) really is an inaccuracy relation, and (b)
show that \(c \succ_{\mathcal{I}, w} c^\prime\) holds for all \(w\),
we'll have shown that \(c\) is accuracy dominated even if we don't have
a numerical measure. That's what I'll make a start on here.

I'll only need to make a start on this project because a sufficient
condition for relative accuracy will be enough for the argument. I'll
leave it as a question for another day whether this kind of approach

Inpsirations

\begin{itemize}
\tightlist
\item
  Kelley (\citeproc{ref-Kelley2026}{2026}); might think this couldn't
  work - but the tempting thought is off by a quantifier shift
\item
  Goodman and Lederman (\citeproc{ref-GoodmanLederemanArXiV}{2024});
  origin of eventually better under any permutation
\item
  Joyce and Weatherson (\citeproc{ref-JoyceWeatherson2019}{2019}); gotta
  do everything with full algebras, not just arbitrary sets of
  propositions
\end{itemize}

We'll now show that if we measure accuracy by the Brier score, if \(A\)
is a merely finitely additive credence function, defined over a Boolean
algebra with a countable infinity of atoms
\(\{w_1, \dots, w_n, \dots\}\) there is some countably additive credence
function \(B\) which accuracy dominates it. The key point is that if
\(A\) is merely finitely additive, then it is possible to define a
countably additive \(B\) such that for each atom
\(w_i, A(\{w_i\}) < B(\{w_i\})\). Since \(A\) is merely finitely
additive, \(\sum_{i \in \mathbb{N}} A(w_i) < 1\). Using this fact, we
can note that if for each \(i\), we define
\(B(\{w_i\}) = A(\{w_i\}) + (1-\sum_{i \in \mathbb{N}} A(w_i))\cdot 2^{-i}\).
If we say \(B\) is countably additive, then for any \(X\),
\(B(X) = \sum_{w_i \in X}B(\{w_i\})\), so saying what values the atoms
take determines the function. Since
\(\sum_{i \in \mathbb{N}}B(\{w_i\}) = 1\), then \(B\) is coherent, and
clearly for each \(w_i, A(\{w_i\}) < B(\{w_i\})\).

Now assume we are using the Brier score to measure inaccuracy, so for
any finite algebra
\(\mathcal{F}, \mathcal{I_F(c, w)} = \sum_{X \in \mathcal{F}}(T(X) - c(X))^2\),
where \(T(X) = 1\) if \(w \in X\) and 0 otherwise.

\begin{theorem}[]\protect\hypertarget{thm-Brier-dom}{}\label{thm-Brier-dom}

If

\begin{itemize}
\tightlist
\item
  \(\mathcal{F}\) is a Boolean algebra whose atoms are some finite
  subset \(\mathcal{F}_0\) of \(\mathbb{N}\),
\item
  the inaccuracy measure \(\mathcal{I_F(c, w)}\) is the Brier score,
\item
  and for any \(X \in \mathcal{F}, A(X) < B(X)\),
\end{itemize}

then for any
\(w \in \mathcal{F}_0, \mathcal{I_F(A, w)} > \mathcal{I_F(B, w)}\)

\end{theorem}

That is, as long as the actual world \(w\) is at atom of the algebra
\(\mathcal{F}\), then \(A\) will be more inaccurate than \(B\) over that
algebra. So no matter how the \(w_i\) are permuted, and whichever world
is actual, eventually there will be a point in the permutated order of
worlds where we reach the actual world, and after that \(B\) will be
more accurate than \(A\). That is, \(B\) will be inevitably eventually
more accurate than \(A\), and so more accurate \emph{simpliciter}.

\begin{proof}[Proof of Theorem~\ref{thm-Brier-dom}]
We can divide \(\mathcal{F}\) into many quadruples of sets of worlds,
each with the following structure. There is some set
\(S \subset \mathcal{F}\) such that \(w \notin S\). Say
\(S^\prime = \mathcal{F} - S - \{w\}\). Then the quadruple will be
\(\langle S, S^\prime, S \cup \{w\}, S^\prime \cup \{w\}\rangle\). The
first two sets of worlds, or as we'll call them from now on,
propositions, will not contain \(w\); they are intuitively false. The
other two are intuitively true. Define the following variables:

\begin{itemize}
\tightlist
\item
  \(x = A(S)\)
\item
  \(x_0 = B(S) - A(S)\), i.e., \(B(S) = x + x_0\)
\item
  \(y = A(S^\prime)\)
\item
  \(y_0 = B(S^\prime) - A(S^\prime)\), i.e., \(B(S^\prime) = y + y_0\)
\item
  \(z = A(\{w\})\)
\item
  \(z_0 = B(\{w\}) - A(\{w\})\), i.e., \(B(\{w\}) = z + z_0\)
\end{itemize}

By construction, each of \(x_0, y_0\) and \(z_0\) are positive. Since
\(S, S^\prime\) and \(\{w\}\) are exclusive, and both \(A\) and \(B\)
are at least finitely additive, it follows that
\(x + x_0 + y + y_0 + z + z_0 \leq 1\), and hence \(x + y + z < 1\).

Now we can work out the Brier score for each of \(A\) and \(B\) over
just the four propositions in this quartet. It helps in the following
argument to remind ourselves of the values of \(A(X)\) and \(B(X)\) for
each of the four propositions in this quartet, as in \$tbl-brier-proof.

{\def\LTcaptype{none} % do not increment counter
\begin{longtable}[]{@{}
  >{\centering\arraybackslash}p{(\linewidth - 6\tabcolsep) * \real{0.1379}}
  >{\centering\arraybackslash}p{(\linewidth - 6\tabcolsep) * \real{0.1724}}
  >{\centering\arraybackslash}p{(\linewidth - 6\tabcolsep) * \real{0.3448}}
  >{\centering\arraybackslash}p{(\linewidth - 6\tabcolsep) * \real{0.3448}}@{}}
\toprule\noalign{}
\begin{minipage}[b]{\linewidth}\centering
Proposition
\end{minipage} & \begin{minipage}[b]{\linewidth}\centering
Truth Value
\end{minipage} & \begin{minipage}[b]{\linewidth}\centering
A
\end{minipage} & \begin{minipage}[b]{\linewidth}\centering
B
\end{minipage} \\
\midrule\noalign{}
\endhead
\bottomrule\noalign{}
\endlastfoot
\(S\) & FALSE & \(x\) & \(x + x_0\) \\
\(S^\prime\) & FALSE & \(y\) & \(y + y_0\) \\
\(S \cup \{w\}\) & TRUE & \(x + z\) & \(x + x_0 + z + z_0\) \\
\(S^\prime\cup \{w\}\) & TRUE & \(y + z\) & \(y + y_0 + z + z_0\) \\
\end{longtable}
}

Say \(\mathcal{I}_4(c, w)\) is the inaccuracy of a credence function
\(c\) over these four propositions. Then we have,

\begin{align*}
\mathcal{I}_4(A, w) &= x^2 + y^2 + (1 - x - z)^2 + (1 - y - z) ^2 \\
&= x^2 + y^2 + 1 + x^2 + z^2 + 2xz - 2x - 2z + \\
&{\hspace{13pt}}1 + y^2 + z^2 + 2yz - 2y - 2z \\
&= 2 + 2x^2 + 2y^2 + 2z^2 + 2xz + 2yz - 2x - 2y - 4z
\end{align*}

\begin{align*}
\mathcal{I}_4(B, w) &= (x+x_0)^2 + (y+y+0)^2 + (1 - x -x_0- z-z_0)^2 + (1 - y -y_0 -z - z_0)^2  \\
&= 7
\end{align*}
\end{proof}

Open Questions

\begin{enumerate}
\def\labelenumi{\arabic{enumi}.}
\tightlist
\item
  Does this generalise to all strictly proper scoring rules?
\item
  Does the converse result, i.e., that countably additive credences are
  not dominated, hold when we use the Brier score?
\item
  Does the converse result hold for all strictly proper scoring rules?
\item
  Do the permutations play any role in the proof, and if not, what role
  do they play in the argument?
\item
  If we bump the Invariably Eventually Better rule from a conditional to
  a biconditional, we violate truth-directedness. (Proof, it treats all
  credences that assign 0 to each atom the same.) Is this a reason to
  not make it a biconditional, or to weaken truth-directedness? (This is
  related to Kelley's discussion of Zhang.)
\item
  How does this all relate to the arguments that are raised in the
  papers in footnote 1.
\end{enumerate}

\protect\phantomsection\label{refs}
\begin{CSLReferences}{1}{0}
\bibitem[\citeproctext]{ref-GoodmanLederemanArXiV}
Goodman, Jeremy, and Harvey Lederman. 2024. {``Maximal Social Welfare
Relations on Infinite Populations Satisfying Permutation Invariance.''}
\url{https://arxiv.org/abs/2408.05851v2}. arXiv preprint.

\bibitem[\citeproctext]{ref-Joyce1998}
Joyce, James M. 1998. {``A Non-Pragmatic Vindication of Probabilism.''}
\emph{Philosophy of Science} 65 (4): 575--603. doi:
\href{https://doi.org/10.1086/392661}{10.1086/392661}.

\bibitem[\citeproctext]{ref-JoyceWeatherson2019}
Joyce, James M., and Brian Weatherson. 2019. {``Accuracy and the
Imps.''} \emph{Logos \& Episteme} 10 (3): 263--82. doi:
\href{https://doi.org/10.5840/logos-episteme201910325}{10.5840/logos-episteme201910325}.

\bibitem[\citeproctext]{ref-Kelley2023}
Kelley, Mikayla. 2023. {``On Accuracy and Coherence with Infinite
Opinion Sets.''} \emph{Philosophy of Science} 90 (1): 92--128. doi:
\href{https://doi.org/10.1017/psa.2021.37}{10.1017/psa.2021.37}.

\bibitem[\citeproctext]{ref-Kelley2026}
---------. 2026. {``A Contextual Accuracy Dominance Argument for
Probabilism.''} \emph{Philosophy and Phenomenological Research} 112 (3):
594--607. doi:
\href{https://doi.org/10.1111/phpr.70098}{10.1111/phpr.70098}.

\bibitem[\citeproctext]{ref-KelleyNeth2023}
Kelley, Mikayla, and Sven Neth. 2023. {``Accuracy and Infinity: A
Dilemma for Subjective {Bayesians}.''} \emph{Synthese} 201 (1): 12. doi:
\href{https://doi.org/10.1007/s11229-022-04019-9}{10.1007/s11229-022-04019-9}.

\bibitem[\citeproctext]{ref-Nielsen2022}
Nielsen, Michael. 2022. {``On the Best Accuracy Arguments for
Probabilism.''} \emph{Philosophy of Science} 89 (3): 621--30.

\bibitem[\citeproctext]{ref-Nielsen2023infinite}
---------. 2023. {``Accuracy and Probabilism in Infinite Domains.''}
\emph{Mind} 132 (526): 402--27. doi:
\href{https://doi.org/10.1093/mind/fzac053}{10.1093/mind/fzac053}.

\bibitem[\citeproctext]{ref-ParkJung2023}
Park, Ilho, and Jaemin Jung. 2023. {``Infinite Opinion Sets and Relative
Accuracy.''} \emph{The Journal of Philosophy} 120 (6): 285--313. doi:
\href{https://doi.org/10.5840/jphil2023120611}{10.5840/jphil2023120611}.

\bibitem[\citeproctext]{ref-Pruss2022}
Pruss, Alexander R. 2022. {``Accuracy, Probabilism and {Bayesian} Update
in Infinite Domains.''} \emph{Synthese} 200 (6): 444. doi:
\href{https://doi.org/10.1007/s11229-022-03938-x}{10.1007/s11229-022-03938-x}.

\bibitem[\citeproctext]{ref-Ross2010}
Ross, Jacob. 2010. {``Sleeping Beauty, Countable Additivity, and
Rational Dilemmas.''} \emph{Philosophical Review} 119 (4): 411--47. doi:
\href{https://doi.org/10.1215/00318108-2010-010}{10.1215/00318108-2010-010}.

\end{CSLReferences}




\end{document}
